%=======================================================================
%  Probing and Steering for Safety-Relevant Behaviours in LLMs
%  Draft v2 — regenerated 2026-08-02 from the v3 pipeline runs.
%
%  Referencing: HARVARD (author--date), per QMUL Faculty of Science and
%  Engineering requirements. Implemented with natbib in author-year mode
%  over a manual thebibliography, so no .bst file is needed:
%     \citep{key} -> (Marks and Tegmark, 2024)
%     \citet{key} -> Marks and Tegmark (2024)
%
%  Build:  latexmk -pdf report_v2.tex
%=======================================================================
\documentclass[conference]{IEEEtran}
\IEEEoverridecommandlockouts

\usepackage[authoryear,round]{natbib}
\usepackage{amsmath,amssymb,amsfonts}
\usepackage{graphicx}
\usepackage{textcomp}
\usepackage{xcolor}
\usepackage{booktabs}
\usepackage{multirow}
\usepackage{url}
\usepackage[hidelinks]{hyperref}

% Harvard-style section heading for the reference list
\renewcommand{\bibsection}{\section*{References}}

% Placeholders for figures not yet exported from the notebooks.
\newcommand{\TODO}[1]{\textcolor{red}{[\textbf{TODO:} #1]}}

\begin{document}

\title{Reading and Writing Safety Behaviours:\\
A Read/Write Dissociation Test for Linear Probes\\
in Llama-3.1-8B-Instruct}

\author{\IEEEauthorblockN{Nimun Kaur Bajwa}
\IEEEauthorblockA{\textit{MSc Artificial Intelligence}, Queen Mary University of London\\
n.k.bajwa@se25.qmul.ac.uk}}

\maketitle

%=======================================================================
\begin{abstract}
Linear activation probes are used both to \emph{detect} safety-relevant
behaviours in language models and, by adding the probe direction back into the
residual stream, to \emph{steer} them. These are different demands on a
direction, and the literature has met them separately: difference-of-means
directions are reported to be the most causally implicated in model outputs,
while logistic-regression probes are reported to be the strongest detectors.
Read together, this suggests a \emph{read/write dissociation}, but no prior work
tests it within a single pipeline on safety behaviours. We build such a pipeline
and apply it uniformly to three behaviours in Llama-3.1-8B-Instruct: deception,
high-stakes awareness, and sycophancy. Each behaviour is specified once and run
through shared machinery: cross-validated layer selection, a five-cell probe
grid crossing estimator with token pooling, in- and out-of-distribution
evaluation by AUROC, and calibrated steering scored by a normalised indirect
effect (NIE) and by an LLM judge that rates behaviour and coherence separately.
The dissociation replicates cleanly on deception and high-stakes: the
logistic-regression family wins 75\% and 100\% of detection columns
respectively, while mass-mean wins 100\% of steering columns in both. It
\emph{fails} on sycophancy, and fails in both directions at once --- mass-mean
is the best detector (0.879 AUROC out-of-distribution) and the worst steerer.
The coherence axis supplies a partial explanation the NIE alone cannot: at high
steering coefficients mass-mean collapses generation quality (judge coherence
falling to 1.0/10 on sycophancy) where logistic-regression and attention
directions do not, so part of mass-mean's apparent causal advantage is
\emph{steering harder}, not steering better. We conclude that the dissociation
is a property of particular behaviour--representation pairs rather than a
general law, and identify the specific measurement conditions under which the
NIE is trustworthy.
\end{abstract}

\begin{IEEEkeywords}
mechanistic interpretability, linear probes, activation steering, AI safety,
deception detection, sycophancy, high-stakes monitoring
\end{IEEEkeywords}

%=======================================================================
\section{Introduction}

As language models are deployed in consequential settings, cheap and reliable
monitoring of their internal state becomes a practical safety requirement.
\emph{Activation probes} --- linear readouts of the residual stream --- are a
leading candidate, flagging a property of interest at a small fraction of the
cost of an auxiliary monitor model \citep{alain2018,mckenzie2025}.

Two distinct questions determine whether such a probe is useful. First, does it
\textbf{read}: does its score separate positive from negative examples, both
in-distribution and, crucially, out-of-distribution? Second, does it
\textbf{write}: is the direction it identifies causally implicated, such that
adding it to the residual stream changes what the model does? A monitoring
system needs the first. An intervention --- suppressing a behaviour at inference
time --- needs the second. The interpretability literature has largely studied
them apart, and on different behaviours.

Three results anchor this work. \citet{marks2024}, studying truth, find that
simple difference-of-means (``mass-mean'') directions are the most causally
implicated, outperforming logistic regression in 7 of 8 intervention conditions,
\emph{despite} all probing techniques classifying at comparable accuracy.
\citet{goldowsky2025}, studying strategic deception, find logistic regression
the strongest detector across diverse evaluation sets. \citet{mckenzie2025},
studying whether an interaction ``might lead to significant harm,'' find that
probes pooling over the whole sequence far outperform last-token probes and
rival monitors orders of magnitude more expensive.

Read together these suggest, but do not test, a \textbf{read/write
dissociation}: the direction that best reads a behaviour off the activations may
not be the direction that best writes it back in. The suggestion is fragile,
because the three results come from different models, different behaviours,
different metrics, and different token-selection policies. Any apparent
dissociation could be an artefact of that heterogeneity rather than a fact about
representations. Testing it requires holding everything else fixed.

\subsection*{Contributions}

\begin{enumerate}
  \item A single specification-driven pipeline in which a behaviour is defined
    once --- data, token mask, readout format, judge rubric --- and inherits
    identical layer selection, probe grid, metrics, and steering machinery.
    Adding a behaviour requires a specification and a loader, nothing else.
  \item The first within-pipeline test of the read/write dissociation on safety
    behaviours, across three behaviours in one model, with in- and
    out-of-distribution evaluation on both sides.
  \item A negative result with diagnostic value: the dissociation replicates on
    deception and high-stakes but inverts on sycophancy, showing it is not a
    general property of linear probes.
  \item Evidence that the coherence axis is necessary, not decorative:
    mass-mean's causal advantage is entangled with generation collapse at high
    coefficients, which a readout-only causal metric cannot detect.
  \item An explicit account of where our protocol diverges from its three source
    papers, and what each divergence costs (Section~\ref{sec:divergences}).
\end{enumerate}

%=======================================================================
\section{Related Work}

\subsection{Linear representations and probing}

The premise underlying this work is the linear representation hypothesis: that
features are encoded as directions in activation space, so a linear readout can
recover them \citep{alain2018,elhage2022}. Probing for truthfulness specifically
was advanced by \citet{azaria2023} and \citet{burns2023}, the latter introducing
Contrast-Consistent Search (CCS), an unsupervised method exploiting the
constraint that a statement and its negation should receive opposite truth
values.

This line attracted immediate criticism. \citet{levinstein2024} showed such
probes fail to generalise in basic ways --- notably to statements containing
``not'' --- suggesting they may track features correlated with truth on the
training distribution rather than truth itself. That critique motivates the
emphasis on out-of-distribution evaluation throughout the present work: an
in-distribution AUROC near 1.0 is, on its own, weak evidence.

\subsection{Geometry of Truth: causality of directions}

\citet{marks2024} address the critique with curated true/false datasets designed
to minimise variation in non-truth features. Three contributions bear directly
on this thesis. First, they localise truth representations by \emph{activation
patching}: running the model on a minimal prompt pair and swapping residual
stream activations at each (position, layer), which isolates a small group of
causally implicated hidden states over the statement-end tokens. Second, they
introduce mass-mean probing, $\boldsymbol{\theta}_{\mathrm{mm}} =
\boldsymbol{\mu}^{+} - \boldsymbol{\mu}^{-}$, together with a
covariance-corrected variant $\sigma(\boldsymbol{\theta}_{\mathrm{mm}}^{\top}
\Sigma^{-1} \mathbf{x})$ used when evaluating in-distribution. Third, and
centrally for us, they quantify causality with a \emph{normalised indirect
effect} (NIE): the shift a direction induces in the model's own TRUE/FALSE
readout, expressed in units of the unsteered gap between classes. Their
headline is that mass-mean directions are the most causal despite comparable
classification accuracy --- the read/write asymmetry this thesis generalises.

\subsection{Deception probes}

\citet{goldowsky2025} train probes on \emph{instructed pairs}: the same fact
presented under honest and dishonest system prompts, so the activation
difference encodes intent rather than content, and the training set contains no
actual deception. Their probe is fitted per-token --- flattening across samples
and sequence positions --- and scored by averaging per-token projections over
the response. They evaluate on realistic settings including AI Liar
\citep{pacchiardi2024}, roleplaying scenarios, and insider trading
\citep{scheurer2023}, reporting AUROCs of 0.96--0.999 and 95--99\% recall at 1\%
false-positive rate, where the false-positive rate is calibrated against a
control set of Alpaca responses \citep{taori2023}. Their appendix benchmarks
difference-of-means and LDA against logistic regression and finds performance
``very similar across methods'' for separating honest from deceptive responses,
with logistic regression winning specifically on separating deception from
\emph{unrelated} transcripts.

\subsection{High-stakes detection}

\citet{mckenzie2025} probe for whether an interaction might lead to significant
harm. Their contribution most relevant here is architectural: they compare
aggregation strategies over the sequence dimension --- mean, max, last token,
softmax-weighted and attention-weighted --- and find learned pooling superior.
They also select their probing layer by 5-fold cross-validation on the training
set, the protocol we adopt.

\subsection{Activation steering}

Contrastive Activation Addition \citep{rimsky2024} computes steering vectors as
mean activation differences over contrast pairs and adds them at inference with
a signed multiplier. Two details matter for our design: CAA adds the vector only
to \emph{generated} tokens, leaving the prompt untouched, and it grades steered
generations with an LLM judge on a 1--10 scale --- the pattern our judge
follows. Related work on inference-time intervention \citep{li2023} and causal
mediation analysis \citep{vig2020,meng2022} supplies the broader framing of
directions as causal mediators.

\subsection{Sycophancy}

Sycophancy --- deferring to a user's stated belief over an accurate answer ---
is well documented as an artefact of preference-based training
\citep{perez2023,sharma2023}. We use the factorial datasets of
\citet{cincynlp2025}, whose design permits a \emph{genuine-agreement} control:
cases where the user happens to be correct, so agreement is candid rather than
sycophantic. This control is what makes sycophancy a discriminating third test
case rather than a repetition of the first two.

\subsection{The gap}

No prior work measures reading and writing for the same behaviour, on the same
model, at the same layers, with matched estimators. That is the gap this thesis
addresses.

%=======================================================================
\section{Methodology}

\subsection{Requirements capture and analysis}

The research question --- \emph{does the direction that best detects a safety
behaviour also best steer it?} --- imposes requirements that are stricter than
those of any single source paper, because the comparison, not the absolute
performance, is the deliverable.

\textbf{R1: Matched arms.} Every probe estimator must be compared under
identical conditions: same layer, same tokens, same calibration, same
intervention site. Any asymmetry becomes a confound.

\textbf{R2: One metric per side, applied uniformly.} The three source papers use
accuracy (\citeauthor{marks2024}), AUROC with recall at 1\% FPR
(\citeauthor{goldowsky2025}) and AUROC (\citeauthor{mckenzie2025}). A single
choice is needed. We use AUROC because our evaluation sets are not
class-balanced and our probe scores are not calibrated to a common threshold ---
notably, mass-mean's raw projection is uncentred --- so threshold-dependent
accuracy would measure calibration drift rather than separability.

\textbf{R3: Out-of-distribution evaluation is mandatory.} Given
\citeauthor{levinstein2024}'s critique, in-distribution performance is not
evidence. Every behaviour requires at least one held-out distribution.

\textbf{R4: Causal claims need a coherence control.} A readout-based causal
metric cannot distinguish ``the direction induced the behaviour'' from ``the
direction broke the model,'' since a degraded model can still emit the target
token. An independent quality measurement is required.

\textbf{R5: Extensibility.} Testing a hypothesis across behaviours requires that
adding the third behaviour not perturb the first two.

\textbf{R6: No selection on evaluation data.} Layers and blocks must be chosen
using training data only, so that out-of-distribution numbers remain
interpretable.

\subsection{Design}

\subsubsection{Behaviour specification}
Each behaviour is one \texttt{BehaviourSpec}: its name, class labels, token mask
policy, the few-shot readout prompt and target tokens for the causal
measurement, the judge rubric, and the steering coefficient sweep. All three
behaviours share the pipeline; only the specification and a data loader differ,
satisfying R5.

\subsubsection{Token masking}
Two policies are used. \texttt{assistant\_only} (deception, sycophancy) probes
the assistant's reply content; \texttt{content\_tokens} (high-stakes) probes
message content of every role. Both exclude BOS, chat-template header tags,
\texttt{<|eot\_id|>} and the template's auto-injected date preamble. Masking is
load-bearing rather than cosmetic: for deception the system prompt \emph{states
the label}, so probing it would read an announcement rather than a
representation.

\subsubsection{Layer selection (R1, R6)}
\label{sec:layers}
Activation patching, \citeauthor{marks2024}'s localisation tool, does not
transfer: it requires single-token minimal pairs and a token-level readout,
whereas our unit is a variable-length dialogue with no canonical
summarisation position. We therefore split the problem by side.

The \emph{detection} layer is the argmax of a pair-aware 5-fold
cross-validated AUROC over mean-pooled training activations, following
\citet{mckenzie2025}. Folds are grouped by pair identifier so that no pair
straddles a fold --- without this, one member of an instructed pair leaks a
near-duplicate of its twin into the test fold. Mass-mean is the selector because
it is closed-form and cheap; a logistic-regression sweep is overlaid as a
robustness check before committing.

The \emph{steering} block is chosen causally: the NIE is computed at every layer
with a single layer steered at a time, and the contiguous five-layer window of
maximum summed NIE is selected. This promotes the NIE from a reporting metric to
a search criterion, which \citeauthor{marks2024} do not do; the consequence is
discussed in Section~\ref{sec:limitations}.

\subsubsection{Probe grid (R1)}
Five cells cross estimator with pooling: mass-mean and logistic regression over
last-token and mean pooling, plus a learned attention probe
\citep{mckenzie2025,kantamneni2025}. The mass-mean probe is
$\boldsymbol{\mu}^{+}-\boldsymbol{\mu}^{-}$. The logistic-regression probe uses
$L_2$ regularisation at $\lambda=10$ on standardised activations, reproducing
\citeauthor{goldowsky2025}'s hyperparameters exactly. The \texttt{LR (mean)}
cell is fitted per-token in their manner; because standardise-then-project is
affine, mean-pooling the activations and projecting equals the mean of
per-token projections, so both their fitting and their aggregation are
reproduced rather than approximated. The flattened token set is uncapped for all
three behaviours: a 60{,}000-row tractability cap inherited from earlier hardware
bound only on high-stakes, where it discarded 79\% of 285{,}394 rows and so made
this one cell a faithful reproduction for two behaviours and a 21\% subsample for
the third. Removing it restores symmetry in the cell whose purpose is to
reproduce \citeauthor{goldowsky2025} exactly.

\subsubsection{Steering (R1, R4)}
Directions are refitted at each block layer and calibrated so that coefficient
$1.0$ moves the mean negative activation onto the mean positive one, following
\citeauthor{marks2024}'s normalisation. This makes coefficients comparable
across estimators whose raw directions differ in norm.

\emph{Where the vector is written} is set by the instrument, and is the one place
this pipeline initially diverged from both source papers without justification.
The vector was originally applied at every token position, a convention inherited
from the ARENA teaching notebook this project builds on \citep{arena2024}. That is
untenable for the readout: the few-shot demonstrations contain the literal
\textsc{yes}/\textsc{no} tokens the metric reads, so steering them can move the
answer logits without changing how the transcript is represented, and the measured
effect would partly be pressure on the instrument.
\citet{marks2024} avoid exactly this by confining the intervention to their
group~(b) hidden states and excluding the readout position. The scope is therefore
narrowed per instrument: the readout steers \textbf{only the item transcript}
(15--21\% of the prompt, varying with transcript length), and generation steers
\textbf{only the completion}, leaving the prompt as the model sees it --- which is
also what real CAA does \citep{rimsky2024}, despite the every-token setting having
cited it. Uniformity across arms is preserved: within a scope every direction
touches identical positions, so a difference between estimators still cannot arise
from one touching fewer tokens than another. \textbf{The steering results in
Section~\ref{sec:results} predate this change and are therefore upper bounds}; see
Section~\ref{sec:limitations}.

Two independent measurements are taken at every (direction, coefficient):

\emph{NIE.} The model reads a few-shot prompt ending in a YES/NO question about
a transcript, and we record the probability difference
$PD = P(\textsc{yes})-P(\textsc{no})$ at the final position. Writing $PD^{+}$
and $PD^{-}$ for the unsteered means over positive and negative transcripts and
$PD^{-}_{*}$ for the steered mean over negatives, the normalised indirect
effect is
\begin{equation}
\mathrm{NIE} \;=\; \frac{PD^{-}_{*} - PD^{-}}{PD^{+} - PD^{-}} ,
\label{eq:nie}
\end{equation}
reproducing \citeauthor{marks2024}'s $\text{false}\!\to\!\text{true}$ ratio with
YES/NO substituted for TRUE/FALSE. $\mathrm{NIE}=0$ means the intervention did
nothing; $\mathrm{NIE}=1$ means it moved the negatives a full class gap. No text
is generated. The denominator makes the metric scale-free, but also fragile: it
is only meaningful when the model can separate the classes unsteered, which
motivates the guard in Section~III-E and the diagnostic in
Table~\ref{tab:gaps}.

\emph{Judge.} The model generates from a prompt with steering active, and Claude
Haiku 4.5 scores the generation 1--10 on the behaviour rubric \emph{and}
separately on coherence, satisfying R4. Scores are bounded server-side by an
enumerated JSON schema.

Only negative-class examples are steered, toward the behaviour; positives supply
the denominator. Coefficients sweep $\{-2,-1,-0.5,0,0.5,1,2\}$, so monotonicity
can be inspected and the negative half acts as a control.

Generation prompts follow one rule across behaviours, adapted from CAA's practice
of deriving open-ended prompts from the same held-out items as the constrained
evaluation \citep{rimsky2024}: \textbf{a generation prompt is one where the
behaviour is off but possible}. For deception that is an AI Liar scenario under its
neutral instruction; for high-stakes, a low-stakes scenario; for sycophancy, a user
turn carrying a \emph{false} claim. The last case needs stating because a user turn
has no sycophancy label of its own --- sycophancy is a property of the reply, which
does not yet exist --- so what a false claim supplies is a prompt on which the
model \emph{could} defer but unsteered does not. A true claim is not
``behaviour off'' but behaviour \emph{impossible}: measured
$P(\text{sycophantic} \mid \text{claim true}) = 0.000$ across 403 rows, structurally,
since sycophancy requires echoing a belief that is wrong. An earlier contiguous
selection of prompts was half true-claim against an 18\% base rate, so half the
sycophancy judge prompts could not register induction at all; selection is now
filtered on claim truth.

\subsection{Implementation}

The pipeline is a Python package (\texttt{safety\_probes}, $\approx$2{,}900
lines) with driver notebooks per behaviour. It evolved through three versions:
an ARENA teaching notebook \citep{arena2024}, a monolithic audited notebook, and
the current modular package, whose docstrings cite the source paper section for
each estimator so provenance is checkable at the line level.

Two implementation decisions were forced by scale. Activations are extracted
once per example at \emph{all} layers in a single forward pass and cached as
fp16, which makes the 32-layer sweep essentially free; full un-pooled token
grids are extracted only at the single chosen layer, since they are roughly
eight times larger. All tokenisation flows through shared helpers, after a
double-BOS defect was found in which templated text and masks disagreed by one
token --- silently scoring every probe on misaligned activations.

Experiments ran on one NVIDIA RTX 5090 (32\,GB) with Llama-3.1-8B-Instruct in
8-bit. The dominant cost is attention-probe training across the steering block:
approximately three hours for deception.

\subsection{Testing and evaluation}

\textbf{Parity gating.} The modular package is required to reproduce the audited
monolithic notebook's selected layers, steering blocks and grid AUROCs before
any result is quoted. Refactors that change numbers are treated as defects.

\textbf{Leakage controls.} Splits are pair-aware throughout: by fact for
deception, by pair for high-stakes, by prompt group for sycophancy. Layer and
block selection use training data only (R6); the block-selection sweep uses
train-split transcripts while reported curves use held-out ones.

\textbf{Metric guards.} AUROC is reported with a bootstrap 95\% confidence
interval (1{,}000 resamples). TPR at 1\% FPR returns \texttt{n/a} when a set has
fewer than 100 negatives, since the target rate is then unreachable and the
metric would silently degrade to TPR at zero false positives --- this is why AI
Liar (27 negatives) is reported as \texttt{n/a} rather than as a number. The NIE
denominator is recomputed on the texts each NIE is measured on, and a warning
fires when the unsteered class gap falls below 0.05, at which point the NIE is
ill-defined.

\textbf{Judge validation.} A blind protocol exists in which the author scores a
sample of generations without seeing the judge's scores, and agreement is
reported as Spearman $\rho$ and mean absolute difference. \emph{This has not yet
been run}, so all judge numbers below are unvalidated; see
Section~\ref{sec:limitations}.

\subsection{Divergences from the source papers}
\label{sec:divergences}

Because the comparison is the deliverable, protocol divergences are stated
rather than absorbed.

\textbf{No covariance-corrected mass-mean arm.} Both \citet{marks2024} (as
$p^{\text{iid}}_{\text{mm}}$) and \citet{goldowsky2025} (as LDA) benchmark
$\Sigma^{-1}(\boldsymbol{\mu}^{+}-\boldsymbol{\mu}^{-})$; we implement only the
plain form. This is well-baselined --- \citeauthor{marks2024} report the
uncorrected probe for every out-of-distribution evaluation and throughout their
causal results --- so the omission costs a possible replication rather than a
missing baseline. It does mean our in-distribution mass-mean numbers should not
be read against \citeauthor{marks2024}'s ``comparable accuracy'' claim, which
rests on the corrected variant.

\textbf{No control set.} \citeauthor{goldowsky2025} calibrate false positives
against Alpaca responses; we lack such a split, so our TPR at 1\% FPR is
computed within each evaluation set. It shares a name with theirs but is a
different measurement, and their headline recall figure is not reproducible
here.

\textbf{Intervention extent.} \citeauthor{marks2024} steer only statement-end
tokens; CAA steers only generated tokens. We steer every token of whatever is
processed, including the readout's few-shot demonstrations. Absolute NIEs are
therefore not comparable with published values.

\textbf{Scale.} All three source papers use 70B-class models; we use 8B, for
which no published baselines exist.

%=======================================================================
\section{Results}
\label{sec:results}

\subsection{Configuration}

Table~\ref{tab:config} reports the selected layers and data budgets. The three
behaviours place their detection layers very differently --- 12, 13 and 19 of 32
--- and their steering blocks differently again, with sycophancy's block
(layers 5--9) lying \emph{below} its detection layer, the only behaviour where
the causal site is earlier than the readable one.

\begin{table}[t]
\caption{Selected configuration and data budgets. Layers are 0-indexed of 32.}
\label{tab:config}
\centering
\small
\begin{tabular}{lccc}
\toprule
& \textbf{Deception} & \textbf{High-stakes} & \textbf{Sycophancy} \\
\midrule
Detection layer & 12 & 13 & 19 \\
Steering block  & 15--19 & 8--12 & 5--9 \\
Mask policy & assistant & content & assistant \\
Train examples & 488 & 3{,}000 & 2{,}736 \\
ID test & 124 & 1{,}000 & 694 \\
OOD set & AI Liar 54 & HH 2{,}984 & math 1{,}000 \\
        & Roleplay 742 & & \\
\bottomrule
\end{tabular}
\end{table}

\begin{figure}[t]
\centering
\fbox{\parbox[c][2.4cm][c]{0.92\columnwidth}{\centering
\TODO{layer-selection panels: pair-aware CV AUROC vs.\ layer for all 32 layers,
MM (selector) and LR (robustness overlay), chosen layer marked. One panel per
behaviour. Source: \texttt{viz.plot\_layer\_selection} in each driver notebook.}}}
\caption{Detection-layer selection by pair-aware 5-fold cross-validated AUROC on
the training split.}
\label{fig:layers}
\end{figure}

\subsection{Linear structure}

For all three behaviours, unsupervised PCA of the detection-layer activations
shows \emph{no} class separation, while the histogram of projections onto the
trained direction separates cleanly. This is the expected and informative
result: the behaviour is a low-variance direction that content variance
swamps. For deception the pair structure is directly visible, with each pair's
two members co-located because they are byte-identical apart from the system
prompt. Had PCA separated the classes, the leading explanation would be a topic
confound of the kind \citeauthor{goldowsky2025} report, where a deception probe
partly tracks ``something close to morality.''

\begin{figure}[t]
\centering
\fbox{\parbox[c][2.4cm][c]{0.92\columnwidth}{\centering
\TODO{PCA + projection panels, one row per behaviour: left, PCA of centred
detection-layer activations (expect no separation); right, histogram of
projections onto the trained MM direction (expect bimodal separation). Source:
\texttt{viz.pca\_and\_projection\_panel}.}}}
\caption{Unsupervised variance does not reveal the behaviour; the supervised
direction does. Left panels show no class separation under PCA; right panels
show clean separation along the probe direction.}
\label{fig:pca}
\end{figure}

\subsection{Detection}

Table~\ref{tab:detection} gives the out-of-distribution AUROC for every cell,
the column on which the read side of the hypothesis is judged.

\begin{table*}[t]
\caption{Detection AUROC with bootstrap 95\% confidence intervals,
in-distribution and out-of-distribution. Deception OOD is the pooled AI Liar
$+$ Roleplaying set; high-stakes OOD is Anthropic-HH; sycophancy OOD is the
held-out math domain. Best OOD cell per behaviour in \textbf{bold}.}
\label{tab:detection}
\centering
\footnotesize
\setlength{\tabcolsep}{3.5pt}
\begin{tabular}{lcccccc}
\toprule
& \multicolumn{2}{c}{\textbf{Deception}} & \multicolumn{2}{c}{\textbf{High-stakes}} & \multicolumn{2}{c}{\textbf{Sycophancy}} \\
\cmidrule(lr){2-3}\cmidrule(lr){4-5}\cmidrule(lr){6-7}
Probe & ID & OOD & ID & OOD & ID & OOD \\
\midrule
MM (last) & 0.948 [0.91,0.98] & 0.696 [0.66,0.73] & 0.867 [0.84,0.89] & 0.678 [0.66,0.70] & 0.976 [0.97,0.98] & \textbf{0.879} [0.86,0.90] \\
MM (mean) & 1.000 [1.00,1.00] & 0.679 [0.64,0.72] & 0.951 [0.94,0.96] & 0.699 [0.68,0.72] & 0.928 [0.91,0.94] & 0.802 [0.77,0.83] \\
LR (last) & 0.996 [0.99,1.00] & 0.571 [0.53,0.61] & 0.977 [0.97,0.98] & 0.782 [0.77,0.80] & 0.975 [0.97,0.98] & 0.847 [0.82,0.87] \\
LR (mean) & 1.000 [1.00,1.00] & \textbf{0.722} [0.69,0.76] & 0.995 [0.99,1.00] & \textbf{0.826} [0.81,0.84] & 0.970 [0.96,0.98] & 0.825 [0.80,0.85] \\
LR-attention & 1.000 [1.00,1.00] & 0.685 [0.65,0.72] & 0.997 [0.99,1.00] & 0.808 [0.79,0.82] & 0.974 [0.96,0.98] & 0.861 [0.84,0.88] \\
\midrule
LR-family win share & \multicolumn{2}{c}{75\%} & \multicolumn{2}{c}{100\%} & \multicolumn{2}{c}{0\%} \\
\bottomrule
\end{tabular}
\end{table*}

Three observations. First, in-distribution AUROC is uninformative: four of five
deception cells reach $\geq 0.996$ and three reach exactly 1.000. The
comparison only becomes visible out-of-distribution, vindicating R3. Second, for
deception and high-stakes the logistic-regression family dominates the OOD
column, as \citeauthor{goldowsky2025} predict --- \texttt{LR (mean)} wins both,
by 0.026 and 0.018 AUROC respectively. Third, sycophancy inverts this entirely:
\texttt{MM (last)} is the best OOD detector at 0.879, and the LR family wins
\emph{no} column.

The deception OOD column decomposes informatively (Table~\ref{tab:decood}). AI
Liar and Roleplaying are deliberately different kinds of held-out set: on AI
Liar the honest and deceptive members share \emph{byte-identical} answer text
and differ only in the system framing, so a probe must read intent in context;
on Roleplaying the completions are genuinely different text, so surface cues are
available. Every cell scores well on AI Liar ($0.818$--$0.992$) and every cell
collapses on Roleplaying ($0.554$--$0.706$).

That ordering is the opposite of what a surface-cue explanation predicts. If the
probes were reading deceptive \emph{phrasing}, Roleplaying --- where the
phrasing differs by label --- should be the easier set. Instead the set where
the text is identical across labels is the one they handle, which is evidence
that the probes read framing-in-context rather than lexical markers. It also
identifies Roleplaying as the genuinely hard generalisation target, and explains
why deception steering fails there entirely (Section~IV-D).

\begin{table}[t]
\caption{Deception detection AUROC split by out-of-distribution set. AI Liar has
byte-identical answers across labels; Roleplaying has genuinely different
completions. TPR at 1\% FPR is \texttt{n/a} on AI Liar (27 negatives $<100$).}
\label{tab:decood}
\centering
\small
\begin{tabular}{lccc}
\toprule
Probe & ID test & AI Liar & Roleplaying \\
\midrule
MM (last)     & 0.948 & 0.959 & 0.669 \\
MM (mean)     & 1.000 & 0.992 & 0.654 \\
LR (last)     & 0.996 & 0.818 & 0.554 \\
LR (mean)     & 1.000 & 0.974 & \textbf{0.706} \\
LR-attention  & 1.000 & 0.989 & 0.657 \\
\bottomrule
\end{tabular}
\end{table}

\subsection{Steering}

Table~\ref{tab:steering} gives the NIE at coefficient $+1$.

\begin{table*}[t]
\caption{Normalised indirect effect at coefficient $+1$. Higher is more causal;
NIE $=1$ means steering moved the negatives one full class gap. Best per column
in \textbf{bold}. The read/write dissociation holds where the LR family wins
Table~\ref{tab:detection} and MM wins here --- true for deception and
high-stakes, false for sycophancy.}
\label{tab:steering}
\centering
\small
\begin{tabular}{lccccccc}
\toprule
& \multicolumn{3}{c}{\textbf{Deception}} & \multicolumn{2}{c}{\textbf{High-stakes}} & \multicolumn{2}{c}{\textbf{Sycophancy}} \\
\cmidrule(lr){2-4}\cmidrule(lr){5-6}\cmidrule(lr){7-8}
Dir. & ID$^{\dagger}$ & Liar & Role & ID & HH & ID & math \\
\midrule
MM  & \textbf{7.25} & \textbf{0.16} & $-$0.07 & \textbf{1.52} & \textbf{0.56} & 0.08 & $-$0.16 \\
LR  & 2.29 & 0.05 & $-$0.08 & 1.01 & 0.12 & $-$0.00 & $-$0.01 \\
LR-att & 1.15 & $-$0.02 & $-$0.18 & 1.38 & 0.26 & \textbf{0.25} & \textbf{0.34} \\
\midrule
MM win share & \multicolumn{3}{c}{100\%} & \multicolumn{2}{c}{100\%} & \multicolumn{2}{c}{0\%} \\
\bottomrule
\end{tabular}
\smallskip
\footnotesize
$^{\dagger}$Unreliable: the unsteered class gap on this set is $+0.058$, barely
above the ill-defined threshold, so the ratio is inflated roughly seventeenfold
(Table~\ref{tab:gaps}). The AI Liar column supports the same ranking on a
healthy gap of $1.052$ and is the one relied on.
\end{table*}

For deception and high-stakes, mass-mean wins every steering column, exactly as
\citeauthor{marks2024} report for truth. Combined with the detection results,
the read/write dissociation \textbf{holds} for both: LR reads best, MM writes
best.

For sycophancy it fails, and fails on both sides at once. Mass-mean is the best
\emph{detector} and among the worst steerers; the attention direction is the
only one producing a clean monotone dose--response (ID: $0.11 \to 0.25 \to
0.51$; math: $0.13 \to 0.34 \to 0.48$). Mass-mean's sycophancy curve is
non-monotone, peaking at coefficient $+0.5$ (ID NIE $0.89$) before collapsing to
$0.08$ at $+1$ --- which is where our verdict is read. The single-point
convention therefore misses a real effect that exists only in a narrow window.

\subsection{When is the NIE trustworthy?}

Because the NIE divides by the unsteered class gap
(Equation~\ref{eq:nie}), its reliability varies by evaluation set in a way the
ratio itself hides. Table~\ref{tab:gaps} reports the denominators.

\begin{table}[t]
\caption{NIE denominators: the unsteered gap $PD^{+}-PD^{-}$ on each evaluation
set. A gap near zero means the model cannot self-report the behaviour in the
readout format, and the ratio becomes unstable.}
\label{tab:gaps}
\centering
\small
\begin{tabular}{llccl}
\toprule
Behaviour & Set & $PD^{+}$ & $PD^{-}$ & Gap \\
\midrule
Deception   & ID          & $-0.906$ & $-0.964$ & $+0.058$ \\
Deception   & AI Liar     & $+0.772$ & $-0.280$ & $+1.052$ \\
Deception   & Roleplaying & $+0.862$ & $+0.122$ & $+0.740$ \\
High-stakes & ID          & $+0.803$ & $+0.463$ & $+0.340$ \\
High-stakes & HH          & $+0.505$ & $-0.894$ & $+1.399$ \\
Sycophancy  & ID          & $+0.982$ & $-0.034$ & $+1.016$ \\
Sycophancy  & math        & $+0.933$ & $+0.162$ & $+0.770$ \\
\bottomrule
\end{tabular}
\end{table}

The deception ID row is the problem case. Both class means sit near $-0.9$: the
model answers ``NO'' to the deception question for honest \emph{and} deceptive
transcripts alike, leaving a gap of $0.058$, barely above the threshold at which
our implementation warns that the metric is ill-defined. The reported ID NIE of
7.25 is that small denominator, not a large effect --- dividing by $0.058$
inflates the ratio roughly seventeenfold relative to a unit gap. The
instructed-pairs format, in which the assistant merely states a fact under a
dishonest persona, apparently does not read as deception to the model judging
its own transcript.

This has a direct methodological consequence: \textbf{the deception ID column
should not be quoted}, and the ranking it appears to support (MM $\gg$ LR $>$
LR-attention) rests on it. The AI Liar column, with a healthy gap of $1.052$,
supports the same ranking on sounder footing and is the one we rely on. The
remaining sets all have gaps above $0.34$ and are usable.

Finally, on Roleplaying all three directions produce NIEs at or below zero:
deception steering does not transfer to the harder OOD set at all, mirroring the
detection collapse on the same set.

\subsection{Coherence: what the NIE cannot see}

Table~\ref{tab:coherence} reports judge coherence at coefficient $+2$ against
the unsteered baseline.

\begin{table}[t]
\caption{Judge coherence (1--10) at coefficient $0$ and $+2$. Mass-mean degrades
generation quality where the other directions do not.}
\label{tab:coherence}
\centering
\small
\begin{tabular}{llcccc}
\toprule
Behaviour & Set & \multicolumn{2}{c}{MM} & LR & LR-att \\
\cmidrule(lr){3-4}
& & coef 0 & coef $+2$ & coef $+2$ & coef $+2$ \\
\midrule
Deception   & Liar & 4.38 & 1.92 & 3.83 & 6.29 \\
Deception   & Role & 6.08 & 1.38 & 2.58 & 4.92 \\
High-stakes & ID   & 4.42 & 2.83 & 4.38 & 5.04 \\
High-stakes & HH   & 6.21 & 2.83 & 5.79 & 6.00 \\
Sycophancy  & ID   & 9.08 & 1.00 & 9.17 & 8.33 \\
Sycophancy  & math & 6.58 & 1.00 & 7.17 & 7.25 \\
\bottomrule
\end{tabular}
\end{table}

The pattern is uniform across all three behaviours: at coefficient $+2$
mass-mean collapses coherence, in the sycophancy case to the floor of the scale,
while logistic-regression and attention directions leave it essentially intact.
On sycophancy this is decisive. Mass-mean does move the judged behaviour score
substantially (ID: $1.17 \to 5.67$ at coefficient $+1$), but coherence falls
from 9.08 to 3.83 at the same point --- the model is producing sycophantic-looking
text by degrading, not by adopting the behaviour cleanly. The attention
direction holds coherence above 8.3 throughout but barely moves the behaviour
score. Neither is a clean success, and the NIE alone would have reported the
first as one.

\begin{figure}[t]
\centering
\fbox{\parbox[c][2.6cm][c]{0.92\columnwidth}{\centering
\TODO{dose--response panels: NIE, judge behaviour and judge coherence against
signed coefficient $(-2 \ldots +2)$, three directions overlaid, one row per
behaviour $\times$ evaluation set. Source:
\texttt{viz.plot\_steering\_effects}. The sycophancy row is the important one:
MM's non-monotone NIE peak at $+0.5$ and its coherence cliff should both be
visible.}}}
\caption{Dose--response. The coherence panel is what distinguishes inducing a
behaviour from degrading the model.}
\label{fig:dose}
\end{figure}

%=======================================================================
\section{Discussion}

\textbf{The dissociation is real but conditional.} On two of three behaviours,
directions that read best and directions that write best are different
directions, reproducing on safety behaviours what \citet{marks2024} found for
truth. This is not an artefact of comparing across papers: here the model,
layer, tokens, calibration and intervention site are identical across arms.

\textbf{Sycophancy shows it is not a law.} The inversion is complete --- the
mass-mean direction is simultaneously the best detector and among the worst
steerers. Several explanations are available and distinguishable in principle.
Sycophancy's label is \emph{relational}: it depends on whether the reply matches
a claim that is itself true or false, and nearly a third of our negatives are
genuine-agreement cases where agreeing is the candid response. A
difference-of-means direction over such data may capture the surface regularity
of agreement, which reads well but is not the causal lever for producing
deference. This is consistent with the observation that sycophancy's steering
block (layers 5--9) sits below its detection layer (19): the site that is
readable is not the site that is writable, which is not the case for the other
two behaviours.

\textbf{Coherence changes the interpretation of causality.} Mass-mean's
advantage on deception and high-stakes co-occurs with systematically worse
generation quality at high coefficients. A readout-based causal metric cannot
see this, because a degraded model can still emit ``YES''. This does not void
the dissociation --- at coefficient $+1$ mass-mean leads on both NIE and judged
behaviour for high-stakes while coherence is close to baseline --- but it means
``more causal'' should be read as ``larger effect per unit coefficient, with
faster degradation,'' not as unambiguously better. This nuance is unavailable to
\citeauthor{marks2024}'s readout-only design and is, we think, the most
transferable methodological finding here.

\textbf{Detection generalisation is the binding constraint.} Every behaviour
loses substantial AUROC out-of-distribution: deception from 1.000 to 0.722,
high-stakes from 0.997 to 0.826, sycophancy from 0.976 to 0.879. At an 8B scale
these probes would not yet support deployment as monitors, which is consistent
with \citeauthor{goldowsky2025}'s own conclusion that current performance is
``insufficient as a robust defence.''

\subsection{Comparison with published results}

Our absolute numbers sit below the source papers', and the gaps are interpretable
rather than mysterious. \citet{goldowsky2025} report AUROCs of 0.96--0.999 on
deception evaluation sets against our 0.722 pooled OOD; \citet{mckenzie2025}
report strong out-of-distribution transfer against our 0.826. Three differences
account for most of this. First, scale: both use 70B-class models and we use 8B,
and \citeauthor{marks2024} show explicitly that probe generalisation improves
with model size. Second, evaluation-set difficulty: our deception OOD pool is
dominated by Roleplaying (742 of 796 examples), which is the harder,
content-defined set, whereas the AI Liar column alone reaches 0.974--0.992 and
is much closer to published territory. Third, metric definition: our TPR at 1\%
FPR is calibrated within each evaluation set rather than against an Alpaca
control, so it is not comparable with \citeauthor{goldowsky2025}'s recall figure
even in principle.

The steering side cannot be compared numerically at all, because our
intervention site differs (every token, versus statement-end tokens in
\citeauthor{marks2024}). What \emph{is} comparable is the ordering, and on
deception and high-stakes the ordering reproduces theirs.

\subsection{Limitations}
\label{sec:limitations}

\textbf{The judge is unvalidated.} The blind agreement protocol has not been
run, so every judge number is provisional. Given that coherence carries much of
the interpretation above, this is the most important outstanding item.

\textbf{NIE is selected on.} Steering blocks are chosen to maximise summed NIE,
and NIE is then reported at those blocks. Train/test separation is maintained
--- the sweep uses training transcripts, reported curves use held-out ones ---
so this is ordinary hyperparameter selection rather than leakage, but
in-distribution NIEs are closer to fitted values than out-of-distribution ones,
and the latter should be given more weight.

\textbf{The block is selected using mass-mean only.} The sweep uses the
mass-mean direction, and the resulting block is used for all three arms. This is
conservative for comparability but leaves an alternative explanation open: if
mass-mean wins on steering, the block may have been chosen to suit it. A sweep
with the logistic-regression direction, showing the block barely moves, would
close this.

\textbf{The reported steering numbers predate the scope fix, and are upper
bounds.} Every NIE in Section~\ref{sec:results} was measured with the vector
written to all tokens of the readout prompt, so part of each value may be direct
pressure on the few-shot demonstrations and the question rather than a change in
how the transcript is represented. The transcript-only scope described in
Section~III-B is implemented and unit-tested but \emph{has not yet been re-run},
and it shrinks the intervention from 100\% of the prompt to 15--21\%. Magnitudes
should be expected to fall. Because the change is uniform across all three
directions, the \emph{ordering} --- which arm writes better, the thesis's actual
claim --- is far more likely to survive than the absolute values. Re-running the
steering sections is the single highest-priority outstanding item, and needs no
re-extraction or probe retraining, since directions and detection results are
unaffected.

\textbf{The Roleplaying judge panel cannot show induction.} Its generation
prompts are scenarios with inherent pressure --- a salesperson who loses
commission by admitting a fault --- and \texttt{roleplaying\_messages} varies only
the assistant completion, so the honest/deceptive flag is a no-op once the
completion is stripped for generation. There is thus no behaviour-off variant of a
Roleplaying prompt, and the unsteered baseline is already 6.96/10 deceptive; the
panel stays within 6.08--7.83 across the whole coefficient sweep for every
direction. That panel should be read as uninformative about induction rather than
as evidence that steering failed. Roleplaying's detection results and its NIE are
unaffected, and its high baseline makes it the natural set for a future
\emph{suppression} experiment.

\textbf{Single model, single seed.} All results are one seed on one 8B model.
The dissociation's dependence on scale is untested.

%=======================================================================
\section{Future Work}

\textbf{Close the open controls.} In priority order: re-run the steering sections
under the corrected transcript-only and completion-only scope, which is
implemented and tested but not yet reflected in any reported number; run the blind
judge validation, without which every judge figure --- including the coherence
result the discussion leans on --- is unvalidated; and sweep the steering block
with the logistic-regression direction to rule out the block having been chosen to
suit mass-mean. A covariance-corrected mass-mean arm was considered and set aside:
both source papers benchmark it, but with $d=4096$ and at most a few thousand
mean-pooled examples the covariance is severely rank-deficient for every behaviour,
so the estimate would depend heavily on the pseudo-inverse regulariser. The cheaper
version of the same test is to report
$\cos(\Sigma^{-1}\boldsymbol{\theta}_{\mathrm{mm}}, \boldsymbol{\theta}_{\mathrm{LR}})$,
which \citeauthor{marks2024}'s App.~F predicts should be near 1 and which would
explain the read-side result mechanistically rather than merely observing it.

\textbf{Complete the metric.} Add the reverse steering direction
($\text{positive}\!\to\!\text{negative}$), which \citeauthor{marks2024} report
and we do not; their results show the two directions can differ substantially.
Add an Alpaca-style control split so that TPR at 1\% FPR means what it means in
the deception literature.

\textbf{Explain the sycophancy inversion.} The relational-label hypothesis is
testable: restricting training to non-genuine-agreement negatives should, if
correct, move the mass-mean direction toward a causal lever at the cost of
detection performance. The read--write gap between layer 19 and layers 5--9
invites a layer-resolved analysis of where deference is computed versus where it
is legible.

\textbf{Scale and mechanism.} Extending to 70B models would test whether the
dissociation strengthens with scale, as \citeauthor{marks2024} observe for
linear structure generally. Replacing linear directions with sparse-autoencoder
latents \citep{cunningham2023,bricken2023} would test whether a sparser basis
separates the readable feature from the writable one --- the natural mechanistic
follow-up to the sycophancy result.

\textbf{Composition.} Preliminary machinery exists for steering two behaviours
simultaneously. Whether directions compose, cancel, or simply produce
incoherence is unanswered by any of the source papers and is a natural extension
given that the coherence axis is already instrumented.

%=======================================================================
\section{Conclusion}

We built a single pipeline that tests, for three safety behaviours in one model,
whether the direction that best detects a behaviour is the direction that best
steers it. For deception and high-stakes the answer is no, and in the specific
manner the literature implies: logistic-regression directions read best,
mass-mean directions write best. For sycophancy the pattern inverts completely,
showing that the dissociation is a property of particular
behaviour--representation pairs rather than of linear probes in general. The
coherence axis further shows that mass-mean's causal advantage is partly a
matter of steering harder rather than steering better. We regard the sycophancy
inversion, and the measurement conditions under which the NIE is and is not
trustworthy, as the more useful contributions of the two.

%=======================================================================
\section*{Acknowledgment}

The author thanks her supervisor for guidance throughout. The pipeline builds on
the ARENA 3.0 interpretability curriculum \citep{arena2024} and on datasets
released by \citet{goldowsky2025}, \citet{mckenzie2025} and
\citet{cincynlp2025}. Use of generative AI tools is declared in Appendix~D.

%=======================================================================
\begin{thebibliography}{99}
\setlength{\itemsep}{2pt}

\bibitem[Alain and Bengio(2018)]{alain2018}
Alain, G. and Bengio, Y. (2018) `Understanding intermediate layers using linear
classifier probes', \textit{arXiv preprint} arXiv:1610.01644.

\bibitem[Azaria and Mitchell(2023)]{azaria2023}
Azaria, A. and Mitchell, T. (2023) `The internal state of an LLM knows when it's
lying', in \textit{Findings of EMNLP 2023}, pp. 967--976.

\bibitem[Bricken et al.(2023)]{bricken2023}
Bricken, T., Templeton, A., Batson, J., et al. (2023) `Towards monosemanticity:
decomposing language models with dictionary learning',
\textit{Transformer Circuits Thread}.

\bibitem[Burns et al.(2023)]{burns2023}
Burns, C., Ye, H., Klein, D. and Steinhardt, J. (2023) `Discovering latent
knowledge in language models without supervision', in \textit{International
Conference on Learning Representations (ICLR)}.

\bibitem[CincyNLP(2025)]{cincynlp2025}
CincyNLP (2025) \textit{Disentangling sycophancy: factorial evaluation
datasets}. Available at: \url{https://github.com/cincynlp/disentangle-sycophancy}
(Accessed: 30 July 2026).

\bibitem[Cunningham et al.(2023)]{cunningham2023}
Cunningham, H., Ewart, A., Riggs, L., Huben, R. and Sharkey, L. (2023) `Sparse
autoencoders find highly interpretable features in language models',
\textit{arXiv preprint} arXiv:2309.08600.

\bibitem[Elhage et al.(2022)]{elhage2022}
Elhage, N., Hume, T., Olsson, C., et al. (2022) `Toy models of superposition',
\textit{Transformer Circuits Thread}.

\bibitem[Goldowsky-Dill et al.(2025)]{goldowsky2025}
Goldowsky-Dill, N., Chughtai, B., Heimersheim, S. and Hobbhahn, M. (2025)
`Detecting strategic deception using linear probes', \textit{arXiv preprint}
arXiv:2502.03407.

\bibitem[Kantamneni et al.(2025)]{kantamneni2025}
Kantamneni, S., Engels, J., Rajamanoharan, S. and Nanda, N. (2025) `Are sparse
autoencoders useful? A case study in sparse probing', \textit{arXiv preprint}
arXiv:2502.16681.

\bibitem[Levinstein and Herrmann(2024)]{levinstein2024}
Levinstein, B.A. and Herrmann, D.A. (2024) `Still no lie detector for language
models: probing empirical and conceptual roadblocks',
\textit{Philosophical Studies}, 181, pp. 2673--2694.

\bibitem[Li et al.(2023)]{li2023}
Li, K., Patel, O., Vi\'egas, F., Pfister, H. and Wattenberg, M. (2023)
`Inference-time intervention: eliciting truthful answers from a language model',
in \textit{Advances in Neural Information Processing Systems (NeurIPS)}.

\bibitem[Marks and Tegmark(2024)]{marks2024}
Marks, S. and Tegmark, M. (2024) `The geometry of truth: emergent linear
structure in large language model representations of true/false datasets', in
\textit{Conference on Language Modeling (COLM)}.

\bibitem[McDougall(2024)]{arena2024}
McDougall, C. (2024) \textit{ARENA 3.0: Alignment Research Engineer
Accelerator}. Available at: \url{https://github.com/callummcdougall/ARENA_3.0}
(Accessed: 2 August 2026).

\bibitem[McKenzie et al.(2025)]{mckenzie2025}
McKenzie, A., Pawar, U., Blandfort, P., Bankes, W., Krueger, D., Lubana, E.S.
and Krasheninnikov, D. (2025) `Detecting high-stakes interactions with
activation probes', \textit{arXiv preprint} arXiv:2506.10805.

\bibitem[Meng et al.(2022)]{meng2022}
Meng, K., Bau, D., Andonian, A. and Belinkov, Y. (2022) `Locating and editing
factual associations in GPT', in \textit{Advances in Neural Information
Processing Systems (NeurIPS)}.

\bibitem[Pacchiardi et al.(2024)]{pacchiardi2024}
Pacchiardi, L., Chan, A.J., Mindermann, S., et al. (2024) `How to catch an AI
liar: lie detection in black-box LLMs by asking unrelated questions', in
\textit{International Conference on Learning Representations (ICLR)}.

\bibitem[Perez et al.(2023)]{perez2023}
Perez, E., Ringer, S., Luko\v{s}i\=ut\.e, K., et al. (2023) `Discovering language
model behaviors with model-written evaluations', in \textit{Findings of ACL
2023}, pp. 13387--13434.

\bibitem[Rimsky et al.(2024)]{rimsky2024}
Rimsky, N., Gabrieli, N., Schulz, J., Tong, M., Hubinger, E. and Turner, A.
(2024) `Steering Llama 2 via contrastive activation addition', in
\textit{Proceedings of ACL 2024}, pp. 15504--15522.

\bibitem[Scheurer et al.(2023)]{scheurer2023}
Scheurer, J., Balesni, M. and Hobbhahn, M. (2023) `Large language models can
strategically deceive their users when put under pressure',
\textit{arXiv preprint} arXiv:2311.07590.

\bibitem[Sharma et al.(2023)]{sharma2023}
Sharma, M., Tong, M., Korbak, T., et al. (2023) `Towards understanding sycophancy
in language models', \textit{arXiv preprint} arXiv:2310.13548.

\bibitem[Taori et al.(2023)]{taori2023}
Taori, R., Gulrajani, I., Zhang, T., et al. (2023) \textit{Stanford Alpaca: an
instruction-following LLaMA model}. Available at:
\url{https://github.com/tatsu-lab/stanford_alpaca} (Accessed: 2 August 2026).

\bibitem[Vig et al.(2020)]{vig2020}
Vig, J., Gehrmann, S., Belinkov, Y., et al. (2020) `Investigating gender bias in
language models using causal mediation analysis', in \textit{Advances in Neural
Information Processing Systems (NeurIPS)}.

\bibitem[Zou et al.(2023)]{zou2023}
Zou, A., Phan, L., Chen, S., et al. (2023) `Representation engineering: a
top-down approach to AI transparency', \textit{arXiv preprint} arXiv:2310.01405.

\end{thebibliography}

%=======================================================================
\appendices

\section{Full Detection Results}
Complete AUROC table with bootstrap 95\% confidence intervals and TPR at 1\%
FPR for all five probe cells $\times$ all evaluation sets, as exported to
\texttt{exported\_probes/*\_detection.csv}.
\TODO{paste the three CSVs as a formatted table; they are already generated.}

\section{Full Dose--Response Curves}
NIE, judge behaviour and judge coherence at each of the seven coefficients, for
every direction and evaluation set, as exported to
\texttt{exported\_probes/*\_steering\_full.json}.
\TODO{paste as three tables.}

\section{Reproducibility}
Model: \texttt{meta-llama/Llama-3.1-8B-Instruct}, 8-bit, seed 42. Hardware:
one NVIDIA RTX 5090 (32\,GB). Selected layers and blocks as
Table~\ref{tab:config}. Data provenance, exact prompts, token masks and
per-example token counts are documented in \texttt{data\_reference.md} and
\texttt{sycophancy\_plan.md} \S5c in the project repository; every prompt in
those files is a real pull from the seed-42 splits.

\section{Statement on the Use of Generative AI Tools}
\TODO{complete using the template provided by the department, and upload
separately as required.}

The following is a factual record to be transferred into that template.
Generative AI tools were used in this project as follows.
\emph{(i) Code.} An AI coding assistant was used to help implement and refactor
the \texttt{safety\_probes} package, to write driver-notebook scaffolding, and
to diagnose defects (for example the double-BOS tokenisation defect described in
Section~III-C). All code was reviewed and executed by the author, and all
reported numbers were produced by running that code.
\emph{(ii) Literature comparison.} An AI assistant was used to cross-check the
implementation against the source papers section by section, producing the
divergence list in Section~\ref{sec:divergences}. The author verified each
divergence against the papers directly.
\emph{(iii) Writing.} An AI assistant was used to draft and edit prose in this
report. The research questions, experimental design, results and conclusions are
the author's own.
\emph{(iv) As an experimental component.} Claude Haiku 4.5 is used
\emph{within} the methodology as the LLM judge scoring steered generations
(Section~III-B); this is a measured component of the system, not a writing aid.

\end{document}
